% 约瑟夫·傅里叶（综述）
% license CCBYSA3
% type Wiki

本文根据 CC-BY-SA 协议转载翻译自维基百科\href{https://en.wikipedia.org/wiki/Joseph_Fourier}{相关文章}

\begin{figure}[ht]
\centering
\includegraphics[width=6cm]{./figures/78d7b0636c3dfd52.png}
\caption{} \label{fig_YSf_1}
\end{figure}
让-巴蒂斯特·约瑟夫·傅里叶（Jean-Baptiste Joseph Fourier，/ˈfʊrieɪ, -iər/；法语：[ʒɑ̃ batist ʒozɛf fuʁje]；1768年3月21日－1830年5月16日）是一位法国数学家和物理学家，出生于勃艮第的欧塞尔。他最为人熟知的是开创了傅里叶级数的研究，这一研究最终发展为傅里叶分析与调和分析，并广泛应用于热传导和振动问题的研究。傅里叶变换和傅里叶导热定律也因他而得名。傅里叶还通常被认为是温室效应的发现者。
\subsection{生平}
傅里叶出生于欧塞尔（今属法国约讷省），父亲是一名裁缝。他在九岁时成为孤儿。后来有人将他推荐给欧塞尔主教，经由这一引荐，他进入圣马克修道院的本笃会接受教育。军队科学部门的任职名额保留给出身高贵的人，而傅里叶因出身寒微不具备资格，因此他接受了一份数学军校讲师的职位。

在法国大革命期间，他在家乡地区积极参与革命活动，并在当地的革命委员会任职。他曾在恐怖统治时期被短暂监禁，但在1795年被任命为师范学校的教师，随后接替约瑟夫-路易·拉格朗日在综合理工学院的职位。

1798年，傅里叶作为拿破仑·波拿巴远征埃及时的科学顾问随行，并被任命为埃及学会的秘书。由于英国舰队的封锁，傅里叶组织了工坊，为法军提供所需的军火。他还为拿破仑在开罗创立的埃及学会（又称开罗学会）撰写了多篇数学论文，该学会旨在削弱英国在东方的影响。1801年，随着英国的胜利以及雅克-弗朗索瓦·梅努率领的法军投降，傅里叶返回法国。
\begin{figure}[ht]
\centering
\includegraphics[width=6cm]{./figures/9ee867e67b3cdc03.png}
\caption{法国艺术家朱利安-莱奥波德·布瓦伊于1820年创作的法国数学家阿德里安-玛丽·勒让德（左）与约瑟夫·傅里叶（右）的水彩漫画，收录于《学会成员的73幅彩色讽刺肖像集》中的第29号和第30号水彩肖像。[3]} \label{fig_YSf_2}
\end{figure}
1801年，[4] 拿破仑任命傅里叶为伊泽尔省（格勒诺布尔）的省长（总督），由他负责监督道路建设及其他工程。然而，在此之前，傅里叶已从拿破仑的埃及远征中返回法国，准备恢复其在综合理工学院的教授职位，但拿破仑在一次评语中作出了不同的决定。

“……由于伊泽尔省的省长近日去世，我愿通过任命傅里叶公民担任此职来表达我对他的信任。[4]”
\begin{figure}[ht]
\centering
\includegraphics[width=6cm]{./figures/ee56ed36e8314992.png}
\caption{克洛德·戈泰罗绘制的傅里叶肖像，约1806年} \label{fig_YSf_3}
\end{figure}
因此，出于对拿破仑的忠诚，他接受了省长一职。[4] 正是在格勒诺布尔任职期间，他开始进行热传导方面的实验。1807年12月21日，他向巴黎学会提交了论文《固体中热的传播》。他还为宏伟的《埃及志》作出了贡献。[5]

1822年，傅里叶接替让-巴蒂斯特·约瑟夫·德朗布尔（Jean Baptiste Joseph Delambre），担任法国科学院的常任秘书。1830年，他当选为瑞典皇家科学院的外国院士。

傅里叶终身未婚。[6]

1830年，他每况愈下的健康终于造成严重影响：

傅里叶早在埃及和格勒诺布尔时就曾遭受过心脏动脉瘤的发作。在巴黎，他频繁出现窒息的主要原因已不可能被误判。然而，1830年5月4日，他在下楼梯时摔了一跤，使得病情加剧到超乎预期的程度。[7]

在这次事故后不久，他于1830年5月16日卧床去世。

傅里叶被安葬在巴黎拉雪兹神父公墓，他的墓碑上装饰着埃及风格的图案，以此体现他曾担任开罗学会秘书以及整理《埃及志》的身份。他的名字也是埃菲尔铁塔上镌刻的72个名字之一。

1849年，欧塞尔为他树立了一座青铜雕像，但在二战期间被熔化制成军械。位于格勒诺布尔的约瑟夫·傅里叶大学则以他的名字命名。
\subsection{《热的解析理论》}
1822年，傅里叶发表了他的著作《热的解析理论》[8]，在书中他将自己的推理建立在牛顿冷却定律的基础上，即：两个相邻质点之间的热流与它们温度的极小差值成正比。该著作在56年后被弗里曼于1878年翻译成英文[9]，并带有编辑性的“修正”[10][11]。此外，数学家让·加斯东·达布鲁对该书进行了大量编辑修订，并于1888年再次以法文出版[10]。

这部著作包含三大重要贡献：一项是纯粹数学的，另外两项则本质上是物理学的。在数学方面，傅里叶声称任意一个变量函数（无论是连续的还是不连续的），都可以展开为该变量整数倍正弦的级数。虽然在没有附加条件时这一结果并不完全正确，但傅里叶关于某些不连续函数可以表示为无穷级数和的观察是一项重大突破。确定傅里叶级数何时收敛的问题，几个世纪以来一直是基础性课题。约瑟夫-路易·拉格朗日曾给出过这一（错误）定理的某些特例，并暗示其方法具有普遍性，但他并未深入研究。后来，彼得·古斯塔夫·勒琼·狄利克雷首次在某些限制条件下给出了令人满意的证明。这项工作为今天被称为傅里叶变换的理论奠定了基础。

在物理学方面，该书包含两个重要贡献：1. 傅里叶提出了方程的量纲齐次性概念，即一个方程形式上只有在等式两边的量纲一致时才是正确的；傅里叶在量纲分析方面作出了重要贡献。[12]2. 他提出了描述热传导扩散的偏微分方程，即著名的热方程。这一方程如今是所有数理物理学生都会学习的内容，也是抛物型偏微分方程最基本的例子。
\subsection{多项式的实根}
\begin{figure}[ht]
\centering
\includegraphics[width=6cm]{./figures/60179e83112ffe0f.png}
\caption{格勒诺布尔的傅里叶半身像} \label{fig_YSf_4}
\end{figure}
傅里叶留下了一部未完成的著作，内容是关于多项式实根的判定与定位，该书由克洛德-路易·纳维耶编辑，并于1831年出版。这部著作包含了许多原创内容——尤其是傅里叶于1820年发表的多项式实根定理。[13][14]在此之前，弗朗索瓦·布当分别于1807年和1811年独立发表了他的定理（也常被称为傅里叶定理），该定理与傅里叶的定理极为接近（两者互为推论）。傅里叶的证明[13] 是19世纪方程论教材中通常采用的标准证明方式。[a]最终，这一问题的完整解答由雅克·夏尔·弗朗索瓦·斯图姆于1829年给出。[15]
\subsection{温室效应的发现}
\begin{figure}[ht]
\centering
\includegraphics[width=6cm]{./figures/919fdcc20ce841a7.png}
\caption{让-巴蒂斯特·约瑟夫·傅里叶之墓，巴黎拉雪兹神父公墓} \label{fig_YSf_6}
\end{figure}
19世纪20年代，傅里叶计算得出：如果地球仅靠入射的太阳辐射加热，那么在其体积和与太阳的距离下，温度应该比实际观测到的要低得多。他在1824年[16]和1827年[17]发表的文章中探讨了额外热量的各种可能来源。然而，由于他的计算结果与实际观测值之间存在高达33摄氏度的巨大差异，傅里叶最终错误地认为星际空间的辐射是重要的热源。尽管如此，傅里叶关于地球大气可能像某种绝热层一样发挥作用的设想，仍被广泛认为是如今所谓温室效应的首次提出[18]，尽管傅里叶本人从未用过这个名称。[19][20]

在文章中，傅里叶引用了奥拉斯-贝内迪克特·德·索叙尔（Horace Bénédict de Saussure）的一项实验。他在一个瓶子内壁铺上了涂黑的软木，并在其中插入几层透明玻璃板，这些玻璃板之间留有空气间隙。正午的阳光透过瓶口的玻璃射入，装置内部越深处的温度越高。傅里叶指出，如果大气中的气体能像这些玻璃板一样形成稳定的屏障，那么它们会对行星温度产生类似的影响。[17]这一结论可能促成了后来“温室效应”这一比喻被用于描述决定大气温度的过程。[21] 傅里叶同时指出，大气温度的实际决定机制还包括对流，而这一点在德·索叙尔的实验装置中并不存在。
\subsection{著作}
\begin{figure}[ht]
\centering
\includegraphics[width=6cm]{./figures/4cd6cb68bf90d7c3.png}
\caption{《热的解析理论》，1888年版} \label{fig_YSf_5}
\end{figure}
\begin{itemize}
\item 《关于笛卡尔定理在寻找根的界限中的应用》。巴黎爱好科学学会公报：156–165，1820年。
\item 《热的解析理论》。巴黎：Firmin Didot Père et Fils，1822年。OCLC 2688081。
\item 《热的解析理论》 ，第1卷。巴黎：Gauthier-Villars，1888年。
\item 《关于地球及行星空间温度的一般性评论》。化学与物理年鉴，27：136–167，1824a。
\item Joseph Louis Gay-Lussac、François Arago 编辑 (1824b)，《辐射热性质的理论综述》。化学与物理年鉴，27，巴黎：236–281。
\item 《关于地球及行星空间温度的论文》。法国皇家科学院院刊，第7卷，1827a，页569–604。W. M. Connolley 英译。
\item 《关于虚根的区分，以及代数分析定理在依赖于热理论的超越方程中的应用》。法国学会皇家科学院院刊 ，第7卷，1827b，页605–624。
\item 《定解方程分析》。Firmin Didot frères，第10卷，1827c，页119–146。2011年9月30日存档。2011年4月20日检索。
\item 《关于代数分析原理在超越方程中的应用的一般性评论》。巴黎：法国学会皇家科学院院刊，第10卷，1827d，页119–146。
\item 《关于流体中热运动的分析论文》。巴黎：法国学会皇家科学院院刊 ，第12卷，1833年，页507–530。
\item 《关于合股保险的报告》。巴黎：法国学会皇家科学院院刊，第5卷，1821年，页26–43。
\end{itemize}
\subsection{参见}
\begin{itemize}
\item 傅里叶分析
\item 傅里叶–德利涅变换 
\item 热方程
\item 最小二乘谱分析 
\item 以约瑟夫·傅里叶命名的事物列表 
\end{itemize}